\section{Symbolic Model and Notation}
\label{sec:notation}

We model the protocol in \mtamarin~\cite{meier2013tamarin}, a symbolic tool for
security-protocol analysis. A system state $S$ is a finite multiset of facts.
Rules are transitions, enabled when their premises are present in the state. A
rule consists of three sets of facts:
\begin{equation}
    [L] \; \text{--}\!\!\text{--}[A] \rightarrow [R]
\end{equation}
where $L$ are the premises, $A$ the action facts, and $R$ the conclusions.
Firing a rule removes the linear premises from the state and adds the
conclusions; action facts are not part of the state but are recorded on the
trace.

\mtamarin\ provides built-in fact symbols \code{In()} and \code{Out()}
representing the adversary-controlled network under
Dolev--Yao~\cite{dolev1983security}, and \code{Fr()} for generating fresh,
unique values. Facts are either \emph{linear} or \emph{persistent}. A linear
fact is consumed when a rule fires; a persistent fact, written with a leading
bang as \code{!F(t)}, is checked but never removed, and may license unboundedly
many rule firings. This distinction is the subject of
\Cref{sec:discipline} and the source of our main modelling finding.

We use the \code{hashing}, \code{signing}, \code{natural-numbers} and
\code{multiset} built-ins; \Cref{sec:nhop} explains why two different
arithmetic builtins are needed. An execution produces a trace of action facts,
and security goals are stated as first-order formulas over traces, proved
either for all traces (safety) or for a single witnessing trace
(exists-trace).
